% Subsistemas del cableado estructurado (ANSI/TIA-568, ANSI/TIA-569)
\begin{tikzpicture}[x=1cm, y=1cm,
  espacio/.style={caja, minimum width=2.5cm, minimum height=0.95cm},
  toma/.style={draw=FigAcento, fill=white, minimum size=0.32cm, inner sep=0pt},
  cable/.style={FigAcento, very thick},
  cordon/.style={FigAcento, thick, densely dotted},
  rot/.style={font=\scriptsize}]
  % Edificio y pisos
  \draw[gray!70, thick] (0,0) rectangle (12,6.6);
  \draw[gray!70] (0,2.2) -- (12,2.2);
  \draw[gray!70] (0,4.4) -- (12,4.4);
  \foreach \y/\piso in {0/Planta baja, 2.2/Piso 1, 4.4/Piso 2}{
    \node[rot, text=gray!70!black, anchor=south east] at (11.95,\y+0.03) {\piso};
  }
  % Planta baja: instalación de entrada y sala de equipos
  \node[espacio] (ef) at (1.6,1.1) {Instalación de\\entrada (EF)};
  \node[espacio] (er) at (4.8,1.1) {Sala de\\equipos (ER)};
  \draw[cable] (-1.7,1.1) -- (ef);
  \node[rot, align=center, above] at (-0.95,1.15) {Proveedor\\externo};
  \draw[cable] (ef) -- (er);
  % Pisos: cuartos de telecomunicaciones y áreas de trabajo
  \foreach \i/\y in {1/3.3, 2/5.5}{
    \node[espacio] (tr\i) at (4.8,\y) {Cuarto de tele-\\comunicaciones (TR)};
    \node[toma] (to\i) at (9.0,\y) {};
    \node[espacio, fill=white, minimum width=1.8cm] (wa\i) at (10.75,\y) {Área de\\trabajo (WA)};
    \draw[cable] (tr\i) -- (to\i);
    \draw[cordon] (to\i) -- (wa\i);
  }
  % Cableado vertebral
  \draw[cable] (er) -- (tr1);
  \draw[cable] (tr1) -- (tr2);
  \node[rot, anchor=east] at (4.65,2.45) {Cableado vertebral};
  % Rótulos y distancias
  \node[rot, below, align=center] at (7.5,5.42) {Cableado\\horizontal};
  \node[rot, above] at (9.0,5.67) {Toma};
  \draw[cota] (6.05,2.75) -- (9.0,2.75) node[midway, below, rot] {Enlace permanente $\leq$ 90 m};
  \draw[cota] (6.05,3.9) -- (11.65,3.9) node[midway, above, rot] {Canal $\leq$ 100 m, incluidos los cordones};
\end{tikzpicture}
